%%
% BIThesis 研究生学位论文模板 The BIThesis Template for Graduate Thesis
% This file has no copyright assigned and is placed in the Public Domain.
%%

\chapter{结论与展望}
\label{chap:conclusion}

随着天地一体化信息网络、低轨卫星星座和在轨智能处理技术的发展，天基网络正在由传统的数据传输网络向集通信、计算、存储和智能处理于一体的分布式协同计算网络演进。在灾害应急通信、遥感快速处理、空间态势感知、区域目标识别和边远区域信息服务等典型场景中，天基网络不仅需要完成数据采集与远程传输，还需要在星地链路受限、星间拓扑高动态、节点资源异构和任务负载突发变化的条件下，实现任务接入、资源协调、计算处理与结果反馈的快速闭环。然而，传统星地集中处理模式容易受到星地链路时延、可见窗口和回传带宽限制，单星本地处理模式又受限于星载平台的算力、能量、功耗和存储资源，难以持续支撑高并发、高时敏和高自治要求的计算任务。因此，如何在高动态、资源受限和去中心化的天基网络环境中，实现可信、高效的多星协同计算卸载，成为提升天基网络智能服务能力的重要问题。

本文围绕高动态天基网络中的可信协同计算卸载问题，按照“统一建模—系统层协同—任务层卸载”的研究主线展开。首先，构建包含动态网络、异构资源、任务生命周期和区块链账本可信状态共享的统一建模框架，并将总体问题分解为系统层协同资源分配和任务层并发开销感知卸载两个子问题；随后，分别设计基于区块链与多智能体强化学习的系统层协同资源分配方法，以及基于聚类联邦学习与DQN的任务层实时卸载方法，从而形成由统一建模支撑、双层方法展开、执行反馈闭环更新的研究体系。全文主要研究工作总结如下。

（1）构建了天基网络可信协同计算卸载统一建模框架。针对天基网络中节点异构、链路时变、任务随机到达和状态共享可信性不足等问题，本文建立了包含卫星终端、卫星边缘节点、地面控制中心和区块链账本的协同计算模型，分别刻画了网络拓扑、节点资源、任务属性、链路传输、本地计算、星间卸载、能耗开销和并发开销等关键要素。在此基础上，本文将区块链账本抽象为可信状态共享与反馈机制，通过链上摘要与链下实体数据分离的方式维护节点状态摘要、任务生命周期记录、节点信誉和模型版本凭据。进一步地，本文建立了统一协同计算卸载优化目标，并将其分解为系统层协同资源分配子问题和任务层并发开销感知卸载子问题，为后续章节提供了统一符号体系、状态接口和问题边界。

（2）提出了基于区块链与多智能体强化学习的系统层协同资源分配方法。针对去中心化天基计算网络中节点能力异构、任务到达非平稳、链路状态动态变化以及节点自利行为导致的资源竞争和负载失衡问题，本文将系统层协同资源分配问题建模为多智能体马尔可夫博弈，并提出区块链增强的集中训练—分散执行多智能体强化学习框架。该方法利用联盟链作为可信公共信息板，维护节点负载、链路状态、信誉评分和历史行为记录等系统级公共摘要，在不依赖长期稳定中心控制器的条件下，为各卫星节点提供可验证、可追溯的协同上下文。同时，本文设计了融合任务完成质量、系统公平性和信誉激励的混合奖励函数，引导自利节点在长期交互中形成更接近系统最优的协同策略。仿真实验表明，该方法能够提升任务完成率和系统吞吐量，改善负载均衡效果，并增强系统对自利行为和不完全信息环境的适应能力。

（3）提出了基于聚类联邦学习与DQN的并发感知任务层计算卸载方法。在系统层实现可信协同的基础上，本文进一步关注任务到达瞬间的节点侧实时卸载决策问题。针对多任务并发执行过程中 CPU 时间片、缓存、内存带宽和 I/O 队列竞争所导致的隐性并发开销，本文建立了任务级并发开销建模方法，将新任务自身计算时延和存量任务剩余时延增加统一纳入任务层代价模型。考虑到天基网络中任务执行日志不适合集中上传、节点数据分布具有明显 Non-IID 特征、区块链账本维护的可信公共摘要存在时滞等问题，本文提出了“聚类式联邦学习预测 + DQN 卸载决策”的闭环机制。该方法通过聚类式联邦学习在不共享原始任务数据的前提下训练并发开销预测模型，并通过区块链账本维护并发开销预测模型版本和可信摘要；随后将预测结果、本地观测状态和区块链账本维护的可信公共摘要共同输入深度 Q 网络，实现低开销、实时化且具备并发开销前瞻能力的任务卸载决策。实验结果表明，该方法能够有效降低高并发场景下的任务时延和能耗，提高任务完成率与系统稳定性。

综上，本文面向天基网络中的可信协同计算卸载需求，形成了由统一建模、系统层协同资源分配和任务层并发开销感知卸载构成的方法体系。研究表明，区块链账本能够为去中心化天基网络提供可信状态共享、任务生命周期记录、信誉维护和并发开销预测模型版本管理能力；多智能体强化学习能够有效支撑系统层多节点协同资源分配；聚类式联邦学习与DQN的结合能够在隐私受限、状态滞后和负载异构条件下提升任务层卸载决策的实时性与准确性。相关研究可为未来构建高效、智能、可信的天基边缘计算网络提供理论依据和方法支撑。

在本文研究基础上，仍有以下方向值得进一步深入探索。

（1）更真实轨道与链路环境下的模型验证。本文主要通过仿真方式验证所提方法，后续可结合真实星座轨道数据、星间链路测量数据和任务负载数据，对模型进行进一步校准。同时，可引入更精细的轨道动力学模型、星间链路中断模型、太阳能充放电模型、热控约束模型和多业务混合到达模型，使统一建模更加贴近真实星座运行环境。

（2）面向大规模星座的轻量化协同优化。随着节点规模扩大，多智能体训练复杂度、区块链账本维护的可信公共摘要维护开销和候选节点筛选成本将进一步增加。后续可研究图神经网络、参数共享、分层策略、模型压缩和边缘推理加速等方式提升可扩展性，使系统能够在更大规模星座中保持较低通信开销和较高协同效率。

（3）区块链可信机制与实时卸载决策的深度融合。本文将区块链主要用于状态摘要、任务生命周期记录、信誉维护和模型版本管理，后续可进一步研究更轻量的共识机制、跨链状态压缩方法、面向时敏任务的可信状态快速确认机制以及链上事件与强化学习奖励之间的自适应耦合方法，从而在可信性和实时性之间取得更优平衡。
